% Complete reviewed Unit 2; editorial master: chapters/ch02/answer-key.md
% Markdown SHA-256: e92678fabfb798725eb0c4d089e6a7e47b20961b402295cc5410ef151033dbab
% One-use initial assembly followed by complete content and page review.
\newcommand{\runningtitle}{Unit 2. Answer key}
% Layout inherited from the reviewed previous course; set a running title per unit.
% This is a small hand-maintained TeX layout, not a conversion pipeline.
\RequirePackage{fix-cm}
\documentclass[11pt,twoside]{article}
\usepackage{fontspec}
\setmainfont{Cambria}[
  Renderer=HarfBuzz,
  BoldFont=cambriab.ttf,
  ItalicFont=cambriai.ttf,
  BoldItalicFont=cambriaz.ttf
]
\setmonofont{lmmono10-regular.otf}[BoldFont=lmmonolt10-bold.otf]
\usepackage[paperwidth=165mm,paperheight=241mm,inner=21mm,outer=19mm,
  top=22mm,bottom=20mm,headheight=14pt,headsep=7mm,footskip=10mm]{geometry}
\usepackage{xcolor,array,booktabs,longtable,tabularx}
\usepackage{enumitem,fancyhdr}
\usepackage[unicode,hidelinks]{hyperref}
\definecolor{ink}{HTML}{000000}
\definecolor{muted}{HTML}{4C5558}
\setlength{\parindent}{0pt}
\setlength{\parskip}{3.5pt}
\AtBeginDocument{\fontsize{11}{14.3}\selectfont}
\setlength{\tabcolsep}{4pt}
\renewcommand{\arraystretch}{1.10}
\setlist[itemize]{leftmargin=1.3em,label=\textbullet,itemsep=2pt,topsep=3pt,parsep=0pt}
\setlist[enumerate]{leftmargin=1.8em,itemsep=3pt,topsep=4pt,parsep=0pt}
\widowpenalty=10000
\clubpenalty=10000
\emergencystretch=1em
\raggedbottom
\pagestyle{fancy}
\fancyhf{}
\providecommand{\runningtitle}{Old English}
\fancyhead[LE,RO]{\small\itshape \runningtitle}
\fancyhead[LO,RE]{\small Old English}
\fancyfoot[LE,RO]{\small\thepage}
\renewcommand{\headrulewidth}{0.3pt}
\newcommand{\booktitle}[2]{
  {\fontsize{9}{11}\selectfont\color{muted}#1\par}
  \vspace{3pt}
  {\fontsize{20}{24}\selectfont\color{ink}#2\par}
  \vspace{6pt}
}
\newcommand{\source}[1]{\par{\footnotesize\color{muted}#1\par}}
\newcommand{\readingtext}[1]{
  \par\addvspace{5pt}
  \begingroup\leftskip=1.5em
    \fontsize{12}{16}\selectfont #1\par
  \endgroup\addvspace{5pt}
}
\makeatletter
\renewcommand{\section}{\@startsection{section}{1}{\z@}
  {-2.4ex \@plus -0.6ex \@minus -0.2ex}{0.8ex \@plus 0.2ex}
  {\normalfont\fontsize{12}{15}\selectfont\bfseries\color{ink}}}
\renewcommand{\subsection}{\@startsection{subsection}{2}{\z@}
  {-1.8ex \@plus -0.5ex \@minus -0.2ex}{0.5ex \@plus 0.2ex}
  {\normalfont\fontsize{11}{14}\selectfont\bfseries\color{ink}}}
\makeatother

\hyphenpenalty=10000
\exhyphenpenalty=10000
\emergencystretch=3em
\newcommand{\itempara}[2]{\par\addvspace{3pt}\noindent\hangindent=1.8em\hangafter=1\makebox[1.8em][l]{#1.}#2\par}
\newcommand{\reserve}[1]{\par\begingroup\dimen0=\pagegoal\advance\dimen0 by-\pagetotal\ifdim\dimen0<#1\newpage\fi\endgroup}
\begin{document}
\raggedright
\booktitle{UNIT 2}{Answer key}

\reserve{7\baselineskip}
% mdparagraph:P001:begin
\section*{Reading the models}
% mdparagraph:P001:end

% mdparagraph:P002:begin
Models use the selected spellings of the lesson. An evidence-supported alternative listed there is valid Old English; where an exercise requests the selected form, distinguish that validity from showing recall of the stated target. Different word order is acceptable when the same cases, agreement, reference and meaning remain clear. Keep the required quantity marks and \textbf{ę}.
% mdparagraph:P002:end

% mdparagraph:P003:begin
A complete verb paradigm includes all three singular and all three plural persons in both tenses. In the compact tables below, the plural row explicitly serves all three persons; dual subjects use that row too. A preterite ending in \textbf{-e} usually has both a first- and a third-singular analysis in these paradigms. Pronouns or context determine which is meant.
% mdparagraph:P003:end

% mdparagraph:P004:begin
Infinitive/\allowbreak{}participle columns are lookup aids, not tests of nonfinite recall or syntax. The six first-/\allowbreak{}second-person genitives and all \textbf{nabban} forms are recognition material. The source-reading help may be consulted throughout its exercises. This key is newly authored; it is not Cook’s own answer key.
% mdparagraph:P004:end

\reserve{7\baselineskip}
% mdparagraph:P005:begin
\section*{9. Pronoun practice}
% mdparagraph:P005:end

\par\begin{minipage}{\linewidth}\raggedright
% item: NEW-U02-P01
% key: KEY-NEW-U02-P01
% mdparagraph:P006:begin
\itempara{1}{Nominative, dative, accusative:}
% mdparagraph:P006:end

% mdtable:T01:begin
\par\begingroup
\fontsize{10.5}{13.2}\selectfont
\setlength{\LTleft}{0pt}\setlength{\LTright}{\fill}
\noindent
\begin{tabular}{@{}>{\raggedright\arraybackslash}p{31mm}>{\raggedright\arraybackslash}p{28mm}>{\raggedright\arraybackslash}p{24mm}>{\raggedright\arraybackslash}p{32mm}@{}}
\toprule
\textbf{Meaning} & \textbf{Nominative} & \textbf{Dative} & \textbf{Accusative} \\
\midrule
I /\allowbreak{} me & ic & mē & mē \\
you, singular & ðū & ðē & ðē \\
\bottomrule
\end{tabular}
\par\endgroup

% mdtable:T01:end
% mdparagraph:P007:begin
The dative may mean “to/\allowbreak{}for me” or “to/\allowbreak{}for you” in a suitable construction; here the giving construction uses it for the recipient. The repeated spellings do not make dative and accusative one case.
% mdparagraph:P007:end

\end{minipage}\par\addvspace{4pt}

\par\begin{minipage}{\linewidth}\raggedright
% item: NEW-U02-P02
% key: KEY-NEW-U02-P02
% mdparagraph:P008:begin
\itempara{2}{\textbf{Wit, unc, unc}, “we two/\allowbreak{}us two”; \textbf{git, inc, inc}, “you two”. The nominatives take \textbf{plural} finite verbs. Retain dual reference in the pronoun; do not invent a dual verbal ending.}
% mdparagraph:P008:end

\end{minipage}\par\addvspace{4pt}

\par\begin{minipage}{\linewidth}\raggedright
% item: NEW-U02-P03
% key: KEY-NEW-U02-P03
% mdparagraph:P009:begin
\itempara{3}{\textbf{Wē, ūs, ūs}; \textbf{gē, ēow, ēow}, in nominative–dative–accusative order. The supplied context asks for ordinary plural forms, rather than the explicitly dual forms.}
% mdparagraph:P009:end

\end{minipage}\par\addvspace{4pt}

\par\begin{minipage}{\linewidth}\raggedright
% item: NEW-U02-P04
% key: KEY-NEW-U02-P04
% mdparagraph:P010:begin
\itempara{4}{Nominative, dative, accusative:}
% mdparagraph:P010:end

% mdtable:T02:begin
\par\begingroup
\fontsize{10.5}{13.2}\selectfont
\setlength{\LTleft}{0pt}\setlength{\LTright}{\fill}
\noindent
\begin{tabular}{@{}>{\raggedright\arraybackslash}p{31mm}>{\raggedright\arraybackslash}p{28mm}>{\raggedright\arraybackslash}p{24mm}>{\raggedright\arraybackslash}p{32mm}@{}}
\toprule
\textbf{Gender} & \textbf{Nominative} & \textbf{Dative} & \textbf{Accusative} \\
\midrule
Masculine & hē & him & hine \\
Feminine & hēo & hiere & hīe \\
Neuter & hit & him & hit \\
\bottomrule
\end{tabular}
\par\endgroup

% mdtable:T02:end
% mdparagraph:P011:begin
Grammatical gender governs the choice. English “it” alone does not determine which row an Old English antecedent requires.
% mdparagraph:P011:end

\end{minipage}\par\addvspace{4pt}

\par\begin{minipage}{\linewidth}\raggedright
% item: NEW-U02-P05
% key: KEY-NEW-U02-P05
% mdparagraph:P012:begin
\itempara{5}{\textbf{Hīe, him, hīe}. No separate plural gender forms are needed. Subject and object \textbf{hīe} share their spelling; their clause roles differ.}
% mdparagraph:P012:end

\end{minipage}\par\addvspace{4pt}

\par\begin{minipage}{\linewidth}\raggedright
% item: NEW-U02-P06
% key: KEY-NEW-U02-P06
% mdparagraph:P013:begin
\itempara{6}{\textbf{Mīn}: first singular; \textbf{ðīn}: second singular; \textbf{uncer}: first dual; \textbf{incer}: second dual; \textbf{ūre}: first plural; \textbf{ēower}: second plural. All are given genitives. This recognition answer does not require composing possessive-adjective phrases.}
% mdparagraph:P013:end

\end{minipage}\par\addvspace{4pt}

\reserve{7\baselineskip}
% mdparagraph:P014:begin
\section*{10. Verb forms}
% mdparagraph:P014:end

\par\begin{minipage}{\linewidth}\raggedright
% item: NEW-U02-F01
% key: KEY-NEW-U02-F01
% mdparagraph:P015:begin
\itempara{1}{\textbf{Hīeran}, “hear”:}
% mdparagraph:P015:end

% mdtable:T03:begin
% course-table-style: 2026-10
% course-panel:T03-P01:begin
\par\begingroup
\fontsize{10.5}{13.2}\selectfont
\setlength{\LTleft}{0pt}\setlength{\LTright}{\fill}
\noindent
\begin{tabular}{@{}>{\raggedright\arraybackslash}p{39.459mm}>{\raggedright\arraybackslash}p{39.459mm}>{\raggedright\arraybackslash}p{39.459mm}@{}}
\toprule
\textbf{Person and number} & \textbf{Present} & \textbf{Preterite} \\
\midrule
1 sg. & hīere & hīerde \\
2 sg. & hīerst & hīerdest \\
3 sg. & hīerð & hīerde \\
1–3 pl. & hīerað & hīerdon \\
\bottomrule
\end{tabular}
\par\endgroup
% course-panel:T03-P01:end
% mdtable:T03:end
% mdparagraph:P016:begin
\textbf{Git hīerdon}, “you two heard”. The dual subject takes the second-person plural form. Present \textbf{hīerst/\allowbreak{}hīerð} and past \textbf{hīerdest/\allowbreak{}hīerde} must remain distinct.
% mdparagraph:P016:end

\end{minipage}\par\addvspace{4pt}

\par\begin{minipage}{\linewidth}\raggedright
% item: NEW-U02-F02
% key: KEY-NEW-U02-F02
% mdparagraph:P017:begin
\itempara{2}{\textbf{Ðū gefylst; hē gefylð; ic gefylde; wē gefyldon.} “You finish; he finishes; I finished; we finished.” Bare correctly labelled verb forms also satisfy the form request. The lexical \textbf{ll} appears in \textbf{gefylle/\allowbreak{}gefyllað}, but simplifies in the selected second/\allowbreak{}third singular before \textbf{st/\allowbreak{}ð} and in the past before \textbf{d}. Do not write an extra \textbf{l} in \textbf{gefylde} or treat the verb here as “fill with a substance”.}
% mdparagraph:P017:end

\end{minipage}\par\addvspace{4pt}

\par\begin{minipage}{\linewidth}\raggedright
% item: NEW-U02-F03
% key: KEY-NEW-U02-F03
% mdparagraph:P018:begin
\itempara{3}{Present 2 singular \textbf{cyst}; present 3 singular \textbf{cyst}; preterite 2 singular \textbf{cystest}; preterite 1/\allowbreak{}2/\allowbreak{}3 plural \textbf{cyston}. The two present cells coincide; their persons do not. The full stem’s \textbf{ss} is not retained in these contracted/\allowbreak{}dental combinations.}
% mdparagraph:P018:end

\end{minipage}\par\addvspace{4pt}

\par\begin{minipage}{\linewidth}\raggedright
% item: NEW-U02-F04
% key: KEY-NEW-U02-F04
% mdparagraph:P019:begin
\itempara{4}{\textbf{Sęttan}: present 1 singular \textbf{sętte}; preterite 1 singular \textbf{sętte}; preterite all plural \textbf{sętton}. \textbf{Sęndan}: present 1 singular \textbf{sęnde}; preterite 1 singular \textbf{sęnde}; preterite all plural \textbf{sęndon}. The identical singular forms need different tense labels. Neither past takes another written dental.}
% mdparagraph:P019:end

\end{minipage}\par\addvspace{4pt}

\par\begin{minipage}{\linewidth}\raggedright
% item: NEW-U02-F05
% key: KEY-NEW-U02-F05
% mdparagraph:P020:begin
\itempara{5}{\textbf{Lǣdan}: present 2 singular \textbf{lǣtst}; present 3 singular \textbf{lǣtt}; preterite 2 singular \textbf{lǣddest}. \textbf{Mētan}: \textbf{mētst, mētt, mēttest} in the same order. Third-singular \textbf{lǣt/\allowbreak{}mēt} are valid final-\textbf{t} alternatives. The long vowels remain long in every requested form; the past of \textbf{lǣdan} has \textbf{dd}, that of \textbf{mētan} \textbf{tt}.}
% mdparagraph:P020:end

\end{minipage}\par\addvspace{4pt}

\par\begin{minipage}{\linewidth}\raggedright
% item: NEW-U02-F06
% key: KEY-NEW-U02-F06
% mdparagraph:P021:begin
\itempara{6}{\textbf{Īecan}: present 1 singular \textbf{īece}; present 3 singular \textbf{īecð}; preterite 1 singular \textbf{īecte}. \textbf{Gierwan}: \textbf{gierwe, giereð, gierede} in the same order. Preserve the selected \textbf{w} in the first form and its absence in the other two. No sentence frame for “increase” or “prepare” is being tested.}
% mdparagraph:P021:end

\end{minipage}\par\addvspace{4pt}

\par\begin{minipage}{\linewidth}\raggedright
% item: NEW-U02-F07
% key: KEY-NEW-U02-F07
% mdparagraph:P022:begin
\itempara{7}{\textbf{Geręstan}, “rest”:}
% mdparagraph:P022:end

% mdtable:T04:begin
% course-table-style: 2026-10
% course-panel:T04-P01:begin
\par\begingroup
\fontsize{10.5}{13.2}\selectfont
\setlength{\LTleft}{0pt}\setlength{\LTright}{\fill}
\noindent
\begin{tabular}{@{}>{\raggedright\arraybackslash}p{39.459mm}>{\raggedright\arraybackslash}p{39.459mm}>{\raggedright\arraybackslash}p{39.459mm}@{}}
\toprule
\textbf{Person and number} & \textbf{Present} & \textbf{Preterite} \\
\midrule
1 sg. & geręste & geręste \\
2 sg. & geręstest & geręstest \\
3 sg. & geręst & geręste \\
1–3 pl. & geręstað & geręston \\
\bottomrule
\end{tabular}
\par\endgroup
% course-panel:T04-P01:end
% mdtable:T04:end
% mdparagraph:P023:begin
\textbf{Hē geręste hine}, “he rested himself/\allowbreak{}he rested”. \textbf{Hē} is nominative, \textbf{hine} accusative and coreferential with it. Contracted present 2 singular \textbf{geręst} is also valid; the table selects \textbf{geręstest}. Past \textbf{geręste} has one \textbf{t}, and its spelling alone permits first or third singular.
% mdparagraph:P023:end

\end{minipage}\par\addvspace{4pt}

\par\begin{minipage}{\linewidth}\raggedright
% item: NEW-U02-F08
% key: KEY-NEW-U02-F08
% mdparagraph:P024:begin
\itempara{8}{\textbf{Sęllan}, “give”:}
% mdparagraph:P024:end

% mdtable:T05:begin
% course-table-style: 2026-10
% course-panel:T05-P01:begin
\par\begingroup
\fontsize{10.5}{13.2}\selectfont
\setlength{\LTleft}{0pt}\setlength{\LTright}{\fill}
\noindent
\begin{tabular}{@{}>{\raggedright\arraybackslash}p{39.459mm}>{\raggedright\arraybackslash}p{39.459mm}>{\raggedright\arraybackslash}p{39.459mm}@{}}
\toprule
\textbf{Person and number} & \textbf{Present} & \textbf{Preterite} \\
\midrule
1 sg. & sęlle & sealde \\
2 sg. & sęlst & sealdest \\
3 sg. & sęlð & sealde \\
1–3 pl. & sęllað & sealdon \\
\bottomrule
\end{tabular}
\par\endgroup
% course-panel:T05-P01:end
% mdtable:T05:end
% mdparagraph:P025:begin
\textbf{Ic sęlle}, “I give”; \textbf{ic sealde}, “I gave”. The changed past stem is lexical. Present \textbf{sęlest/\allowbreak{}sęleð} are valid alternatives to \textbf{sęlst/\allowbreak{}sęlð}; they do not produce a regular past in \textbf{-ede}.
% mdparagraph:P025:end

\end{minipage}\par\addvspace{4pt}

\par\begin{minipage}{\linewidth}\raggedright
% item: NEW-U02-F09
% key: KEY-NEW-U02-F09
% mdparagraph:P026:begin
\itempara{9}{Present 2 singular \textbf{gewyrcst}; present 3 singular \textbf{gewyrcð}. Preterite: 1 singular \textbf{geworhte}, 2 singular \textbf{geworhtest}, 3 singular \textbf{geworhte}. Retain \textbf{ge-} and the past \textbf{worht-} stem. This is the paradigm of \textbf{gewyrcean}, not a general transformation licensed for every prefixed verb.}
% mdparagraph:P026:end

\end{minipage}\par\addvspace{4pt}

\par\begin{minipage}{\linewidth}\raggedright
% item: NEW-U02-F10
% key: KEY-NEW-U02-F10
% mdparagraph:P027:begin
\itempara{10}{Present: 1 singular \textbf{gefręmme}, 2 singular \textbf{gefręmest}, 3 singular \textbf{gefręmeð}, all plural \textbf{gefręmmað}. Preterite 1 singular \textbf{gefręmede}. The \textbf{mm/\allowbreak{}m} alternation belongs to this short-stem family; neither a uniform doubled consonant nor the \textbf{gefylde} pattern gives the selected answers.}
% mdparagraph:P027:end

\end{minipage}\par\addvspace{4pt}

\par\begin{minipage}{\linewidth}\raggedright
% item: NEW-U02-F11
% key: KEY-NEW-U02-F11
% mdparagraph:P028:begin
\itempara{11}{Present 1 singular \textbf{ęfne}; present 2 singular \textbf{ęfnest}; preterite 1 singular \textbf{ęfnede}; preterite all plural \textbf{ęfnedon}. Source-supported \textbf{ęfnde/\allowbreak{}ęfndon} are valid contracted past alternatives. This task tests the forms, not an independently assigned complement construction.}
% mdparagraph:P028:end

\end{minipage}\par\addvspace{4pt}

\par\begin{minipage}{\linewidth}\raggedright
% item: NEW-U02-F12
% key: KEY-NEW-U02-F12
% mdparagraph:P029:begin
\itempara{12}{\textbf{Nęrian}, “save”:}
% mdparagraph:P029:end

% mdtable:T06:begin
% course-table-style: 2026-10
% course-panel:T06-P01:begin
\par\begingroup
\fontsize{10.5}{13.2}\selectfont
\setlength{\LTleft}{0pt}\setlength{\LTright}{\fill}
\noindent
\begin{tabular}{@{}>{\raggedright\arraybackslash}p{39.459mm}>{\raggedright\arraybackslash}p{39.459mm}>{\raggedright\arraybackslash}p{39.459mm}@{}}
\toprule
\textbf{Person and number} & \textbf{Present} & \textbf{Preterite} \\
\midrule
1 sg. & nęrie & nęrede \\
2 sg. & nęrest & nęredest \\
3 sg. & nęreð & nęrede \\
1–3 pl. & nęriað & nęredon \\
\bottomrule
\end{tabular}
\par\endgroup
% course-panel:T06-P01:end
% mdtable:T06:end
% mdparagraph:P030:begin
\textbf{Wit nęriað}, “we two save”; \textbf{ðū nęredest}, “you saved”. Retain present \textbf{i} in the first singular and plural, but not in \textbf{nęrest/\allowbreak{}nęreð} or the selected past \textbf{nęred-}. Dual \textbf{wit} takes the plural verb.
% mdparagraph:P030:end

\end{minipage}\par\addvspace{4pt}

\par\begin{minipage}{\linewidth}\raggedright
% item: NEW-U02-F13
% key: KEY-NEW-U02-F13
% mdparagraph:P031:begin
\itempara{13}{\textbf{Lufian}, “love”:}
% mdparagraph:P031:end

% mdtable:T07:begin
% course-table-style: 2026-10
% course-panel:T07-P01:begin
\par\begingroup
\fontsize{10.5}{13.2}\selectfont
\setlength{\LTleft}{0pt}\setlength{\LTright}{\fill}
\noindent
\begin{tabular}{@{}>{\raggedright\arraybackslash}p{39.459mm}>{\raggedright\arraybackslash}p{39.459mm}>{\raggedright\arraybackslash}p{39.459mm}@{}}
\toprule
\textbf{Person and number} & \textbf{Present} & \textbf{Preterite} \\
\midrule
1 sg. & lufie & lufode \\
2 sg. & lufast & lufodest \\
3 sg. & lufað & lufode \\
1–3 pl. & lufiað & lufodon \\
\bottomrule
\end{tabular}
\par\endgroup
% course-panel:T07-P01:end
% mdtable:T07:end
% mdparagraph:P032:begin
The present second-singular ending contains \textbf{a}, \textbf{-ast}; the preterite has \textbf{o}, \textbf{-od-}. Its infinitive \textbf{-ian} does not justify Class I \textbf{-est/\allowbreak{}-ede}. \textbf{Lufedon} is a valid source plural alternative; \textbf{lufodon} is the selected target.
% mdparagraph:P032:end

\end{minipage}\par\addvspace{4pt}

\par\begin{minipage}{\linewidth}\raggedright
% item: NEW-U02-F14
% key: KEY-NEW-U02-F14
% mdparagraph:P033:begin
\itempara{14}{\textbf{Geblētsian}: present 3 singular \textbf{geblētsað}; preterite 2 singular \textbf{geblētsodest}; preterite all plural \textbf{geblētsodon}. \textbf{Gehālgian}: \textbf{gehālgað, gehālgodest, gehālgodon} in the same order. The prefixes and the lexical quantities \textbf{ē/\allowbreak{}ā} remain; \textbf{-að} is singular here, whereas \textbf{-iað} is the selected present plural.}
% mdparagraph:P033:end

\end{minipage}\par\addvspace{4pt}

\par\begin{minipage}{\linewidth}\raggedright
% item: NEW-U02-F15
% key: KEY-NEW-U02-F15
% mdparagraph:P034:begin
\itempara{15}{\textbf{Habban}, “have, possess”:}
% mdparagraph:P034:end

% mdtable:T08:begin
% course-table-style: 2026-10
% course-panel:T08-P01:begin
\par\begingroup
\fontsize{10.5}{13.2}\selectfont
\setlength{\LTleft}{0pt}\setlength{\LTright}{\fill}
\noindent
\begin{tabular}{@{}>{\raggedright\arraybackslash}p{39.459mm}>{\raggedright\arraybackslash}p{39.459mm}>{\raggedright\arraybackslash}p{39.459mm}@{}}
\toprule
\textbf{Person and number} & \textbf{Present} & \textbf{Preterite} \\
\midrule
1 sg.: ic & hæbbe & hæfde \\
2 sg.: ðū & hæfst & hæfdest \\
3 sg.: hē /\allowbreak{} hēo /\allowbreak{} hit & hæfð & hæfde \\
1–3 pl.: wē /\allowbreak{} gē /\allowbreak{} hīe & habbað & hæfdon \\
\bottomrule
\end{tabular}
\par\endgroup
% course-panel:T08-P01:end
% mdtable:T08:end
% mdparagraph:P035:begin
Each slash-separated pronoun has the verb printed in its row. The third singular may use any appropriate \textbf{hē/\allowbreak{}hēo/\allowbreak{}hit}; the plural row represents each of first, second and third person. Source variants \textbf{hafast, hafað, hæbbað} remain valid. The requested paradigm is affirmative; \textbf{næbbe} is a negative form.
% mdparagraph:P035:end

\end{minipage}\par\addvspace{4pt}

\par\begin{minipage}{\linewidth}\raggedright
% item: NEW-U02-F16
% key: KEY-NEW-U02-F16
% mdparagraph:P036:begin
\itempara{16}{\textbf{Libban}, “live”:}
% mdparagraph:P036:end

% mdtable:T09:begin
% course-table-style: 2026-10
% course-panel:T09-P01:begin
\par\begingroup
\fontsize{10.5}{13.2}\selectfont
\setlength{\LTleft}{0pt}\setlength{\LTright}{\fill}
\noindent
\begin{tabular}{@{}>{\raggedright\arraybackslash}p{39.459mm}>{\raggedright\arraybackslash}p{39.459mm}>{\raggedright\arraybackslash}p{39.459mm}@{}}
\toprule
\textbf{Person and number} & \textbf{Present} & \textbf{Preterite} \\
\midrule
1 sg. & libbe & lifde \\
2 sg. & leofast & lifdest \\
3 sg. & leofað & lifde \\
1–3 pl. & libbað & lifdon \\
\bottomrule
\end{tabular}
\par\endgroup
% course-panel:T09-P01:end
% mdtable:T09:end
% mdparagraph:P037:begin
Keep \textbf{libb-} in the selected first singular/\allowbreak{}plural present, \textbf{leof-} in the other present singular persons, and \textbf{lifd-} in the past. \textbf{Lifiað} is a valid alternative present plural; the selected table is not a regular \textbf{lufian} paradigm.
% mdparagraph:P037:end

\end{minipage}\par\addvspace{4pt}

\par\begin{minipage}{\linewidth}\raggedright
% item: NEW-U02-F17
% key: KEY-NEW-U02-F17
% mdparagraph:P038:begin
\itempara{17}{\textbf{Sęcgan}, “say”:}
% mdparagraph:P038:end

% mdtable:T10:begin
% course-table-style: 2026-10
% course-panel:T10-P01:begin
\par\begingroup
\fontsize{10.5}{13.2}\selectfont
\setlength{\LTleft}{0pt}\setlength{\LTright}{\fill}
\noindent
\begin{tabular}{@{}>{\raggedright\arraybackslash}p{39.459mm}>{\raggedright\arraybackslash}p{39.459mm}>{\raggedright\arraybackslash}p{39.459mm}@{}}
\toprule
\textbf{Person and number} & \textbf{Present} & \textbf{Preterite} \\
\midrule
1 sg. & sęcge & sægde \\
2 sg. & sægst & sægdest \\
3 sg. & sægð & sægde \\
1–3 pl. & sęcgað & sægdon \\
\bottomrule
\end{tabular}
\par\endgroup
% course-panel:T10-P01:end
% mdtable:T10:end
% mdparagraph:P039:begin
Documented alternatives include \textbf{sęgst/\allowbreak{}sagast}, \textbf{sęgð/\allowbreak{}sagað}, plural \textbf{sęcgeað}, and the past series \textbf{sǣde, sǣdest, sǣdon}. Use a source-supported series and preserve its quantities. No indirect-speech construction is being tested.
% mdparagraph:P039:end

\end{minipage}\par\addvspace{4pt}

\par\begin{minipage}{\linewidth}\raggedright
% item: NEW-U02-F18
% key: KEY-NEW-U02-F18
% mdparagraph:P040:begin
\itempara{18}{\textbf{Hycgan}, “think”:}
% mdparagraph:P040:end

% mdtable:T11:begin
% course-table-style: 2026-10
% course-panel:T11-P01:begin
\par\begingroup
\fontsize{10.5}{13.2}\selectfont
\setlength{\LTleft}{0pt}\setlength{\LTright}{\fill}
\noindent
\begin{tabular}{@{}>{\raggedright\arraybackslash}p{39.459mm}>{\raggedright\arraybackslash}p{39.459mm}>{\raggedright\arraybackslash}p{39.459mm}@{}}
\toprule
\textbf{Person and number} & \textbf{Present} & \textbf{Preterite} \\
\midrule
1 sg. & hycge & hogode \\
2 sg. & hygst & hogodest \\
3 sg. & hygð & hogode \\
1–3 pl. & hycgað & hogodon \\
\bottomrule
\end{tabular}
\par\endgroup
% course-panel:T11-P01:end
% mdtable:T11:end
% mdparagraph:P041:begin
Documented alternatives include \textbf{hogast/\allowbreak{}hogað}, plural \textbf{hycgeað}, and past \textbf{hogde, hogdest, hogdon}. The full selected morphology does not itself provide an indirect question or other complement frame.
% mdparagraph:P041:end

\end{minipage}\par\addvspace{4pt}

\par\begin{minipage}{\linewidth}\raggedright
% item: NEW-U02-F19
% key: KEY-NEW-U02-F19
% mdparagraph:P042:begin
\itempara{19}{All seven given forms belong to negative \textbf{nabban}:}
% mdparagraph:P042:end

% mdtable:T12:begin
\par\begingroup
\fontsize{10.5}{13.2}\selectfont
\setlength{\LTleft}{0pt}\setlength{\LTright}{\fill}
\noindent
\begin{tabular}{@{}>{\raggedright\arraybackslash}p{24mm}>{\raggedright\arraybackslash}p{49mm}>{\raggedright\arraybackslash}p{44mm}@{}}
\toprule
\textbf{Form} & \textbf{Analysis} & \textbf{Meaning} \\
\midrule
næbbe & present 1 singular & I do not have \\
næfst & present 2 singular & you do not have \\
næfð & present 3 singular & he/\allowbreak{}she/\allowbreak{}it does not have \\
nabbað & present 1, 2 or 3 plural & we/\allowbreak{}you/\allowbreak{}they do not have \\
næfde & preterite 1 or 3 singular & I/\allowbreak{}he/\allowbreak{}she/\allowbreak{}it did not have \\
næfdest & preterite 2 singular & you did not have \\
næfdon & preterite 1, 2 or 3 plural & we/\allowbreak{}you/\allowbreak{}they did not have \\
\bottomrule
\end{tabular}
\par\endgroup

% mdtable:T12:end
% mdparagraph:P043:begin
The plural verb forms also occur with appropriate dual subjects. The given spelling alone does not select a plural person or distinguish the two singular persons of \textbf{næfde}. Recognition is sufficient; unaided production was not requested.
% mdparagraph:P043:end

\end{minipage}\par\addvspace{4pt}

\par\begin{minipage}{\linewidth}\raggedright
% item: NEW-U02-F20
% key: KEY-NEW-U02-F20
% mdparagraph:P044:begin
\itempara{20}{\textbf{Hīerde → hīeran}, “heard”; \textbf{sealde → sęllan}, “gave”; \textbf{geworhte → gewyrcean}, “made”; \textbf{gefręmede → gefręmman}, “performed/\allowbreak{}accomplished”; \textbf{lufode → lufian}, “loved”; \textbf{lifde → libban}, “lived”. All six are preterites and all permit \textbf{first or third singular indicative}. The infinitive citation identifies the lemma; it is not the form to insert into the sentence.}
% mdparagraph:P044:end

\end{minipage}\par\addvspace{4pt}

\reserve{7\baselineskip}
% mdparagraph:P045:begin
\section*{11. Cases, reference and clauses}
% mdparagraph:P045:end

\par\begin{minipage}{\linewidth}\raggedright
% item: NEW-U02-C01
% key: KEY-NEW-U02-C01
% mdparagraph:P046:begin
\itempara{1}{\textbf{“He hears them. They hear him.”} In the first clause \textbf{hē} is nominative masculine singular, the man; \textbf{hīe} is accusative plural, the birds. In the second \textbf{hīe} is nominative plural, the birds; \textbf{hine} is accusative masculine singular, the man. The hearer and the entity heard exchange roles, so \textbf{hīerð} becomes plural \textbf{hīerað}. The specified referents rule out an unrelated feminine singular reading of the first object \textbf{hīe}.}
% mdparagraph:P046:end

\end{minipage}\par\addvspace{4pt}

\par\begin{minipage}{\linewidth}\raggedright
% item: NEW-U02-C02
% key: KEY-NEW-U02-C02
% mdparagraph:P047:begin
\itempara{2}{\textbf{Ic sęlle ðǣm męnn sǣd. Ic sęlle him sǣd.} The man is the dative recipient; \textbf{him} replaces that entire noun phrase. \textbf{Sǣd} is the accusative thing given. Accusative \textbf{hine} would not preserve the recipient construction; the noun’s dative singular is \textbf{męnn}.}
% mdparagraph:P047:end

\end{minipage}\par\addvspace{4pt}

\par\begin{minipage}{\linewidth}\raggedright
% item: NEW-U02-C03
% key: KEY-NEW-U02-C03
% mdparagraph:P048:begin
\itempara{3}{\textbf{“We two have food. The man gives us two water.”} \textbf{Wit} is nominative dual, subject of plural \textbf{habbað}. \textbf{Unc} is dative dual, the recipient of \textbf{wæter}. It does not control singular \textbf{sęlð}; nominative \textbf{se mann} does.}
% mdparagraph:P048:end

\end{minipage}\par\addvspace{4pt}

\par\begin{minipage}{\linewidth}\raggedright
% item: NEW-U02-C04
% key: KEY-NEW-U02-C04
% mdparagraph:P049:begin
\itempara{4}{\textbf{Git nęriað mē. Ic hīere inc.} \textbf{Git} is nominative dual and takes plural \textbf{nęriað}; \textbf{mē} is accusative singular. In the second clause \textbf{inc} is accusative dual, the two people addressed. Singular \textbf{hīere} agrees with \textbf{ic}.}
% mdparagraph:P049:end

\end{minipage}\par\addvspace{4pt}

\par\begin{minipage}{\linewidth}\raggedright
% item: NEW-U02-C05
% key: KEY-NEW-U02-C05
% mdparagraph:P050:begin
\itempara{5}{\textbf{Sēo eorðe hæfð wæter. Hēo hæfð gærs.} “The earth has water. It has grass.” \textbf{Eorðe} is feminine, so the replacement nominative is \textbf{hēo}. English “it” is natural; Old English \textbf{hit} would not preserve the stated noun’s gender.}
% mdparagraph:P050:end

\end{minipage}\par\addvspace{4pt}

\par\begin{minipage}{\linewidth}\raggedright
% item: NEW-U02-C06
% key: KEY-NEW-U02-C06
% mdparagraph:P051:begin
\itempara{6}{\textbf{His hof is on ðǣre eorðan. His fuglas libbað on ðǣm hofe.} The male possessor gives \textbf{his} in both clauses. Singular \textbf{hof} and plural \textbf{fuglas} determine \textbf{is/\allowbreak{}libbað}, not the possessor. Both locations take dative: \textbf{eorðan} and \textbf{hofe}.}
% mdparagraph:P051:end

\end{minipage}\par\addvspace{4pt}

\par\begin{minipage}{\linewidth}\raggedright
% item: NEW-U02-C07
% key: KEY-NEW-U02-C07
% mdparagraph:P052:begin
\itempara{7}{\textbf{Hēo hæfð hof. Fuglas libbað on hiere hofe.} \textbf{Hiere} identifies the female possessor; neuter dative \textbf{hofe} does not change it to \textbf{his/\allowbreak{}him}. The noun takes dative because it gives the location after \textbf{on}.}
% mdparagraph:P052:end

\end{minipage}\par\addvspace{4pt}

\par\begin{minipage}{\linewidth}\raggedright
% item: NEW-U02-C08
% key: KEY-NEW-U02-C08
% mdparagraph:P053:begin
\itempara{8}{\textbf{Hīe habbað hof. Hiera hof hæfð fuglas.} The first subject is plural \textbf{hīe}. The second subject is singular \textbf{hof}, possessed by the plural group; therefore its verb is \textbf{hæfð}. \textbf{Hiera} remains the invariant genitive possessor.}
% mdparagraph:P053:end

\end{minipage}\par\addvspace{4pt}

\par\begin{minipage}{\linewidth}\raggedright
% item: NEW-U02-C09
% key: KEY-NEW-U02-C09
% mdparagraph:P054:begin
\itempara{9}{\textbf{Ic geręste mē. Hē geręste hine.} Subjects: nominative \textbf{ic}, \textbf{hē}. Objects: accusative \textbf{mē}, \textbf{hine}, each referring back to its subject. The past verb is identical in first and third singular; that does not make the subjects interchangeable. Natural English may omit the reflexive after “rested”.}
% mdparagraph:P054:end

\end{minipage}\par\addvspace{4pt}

\par\begin{minipage}{\linewidth}\raggedright
% item: NEW-U02-C10
% key: KEY-NEW-U02-C10
% mdparagraph:P055:begin
\itempara{10}{\textbf{Git geręstað inc. Hīe geręstað hīe.} The first object is accusative dual \textbf{inc}; the second is accusative plural \textbf{hīe}. Both verbs use the plural form, but the pronouns preserve the different number references. In the second clause the first \textbf{hīe} is nominative, the second accusative.}
% mdparagraph:P055:end

\end{minipage}\par\addvspace{4pt}

\par\begin{minipage}{\linewidth}\raggedright
% item: NEW-U02-C11
% key: KEY-NEW-U02-C11
% mdparagraph:P056:begin
\itempara{11}{\textbf{Ic ne hīere ðone fugol. Hē ne lufað ðā fuglas.} “I do not hear the bird. He does not love the birds.” \textbf{Ne} precedes each finite verb. English “do/\allowbreak{}does” is part of the translation, not an extra Old English verb.}
% mdparagraph:P056:end

\end{minipage}\par\addvspace{4pt}

\par\begin{minipage}{\linewidth}\raggedright
% item: NEW-U02-C12
% key: KEY-NEW-U02-C12
% mdparagraph:P057:begin
\itempara{12}{\textbf{Ðū ne hīerdest hine. Wē ne sęndon ðā fuglas.} Keep past second-singular \textbf{-dest} and past plural \textbf{-on}. \textbf{Hine} and \textbf{ðā fuglas} are accusative objects. Present \textbf{hīerst/\allowbreak{}sęndað} would change the tense.}
% mdparagraph:P057:end

\end{minipage}\par\addvspace{4pt}

\par\begin{minipage}{\linewidth}\raggedright
% item: NEW-U02-C13
% key: KEY-NEW-U02-C13
% mdparagraph:P058:begin
\itempara{13}{\textbf{Ic sęlle mē męte. Ic sęlle him męte.} Both \textbf{mē} and \textbf{him} are dative recipients; \textbf{męte} is the accusative thing given. Only \textbf{mē} is reflexive, since it refers to the subject \textbf{ic}. The dative form \textbf{him} does not mean “himself” merely by its spelling.}
% mdparagraph:P058:end

\end{minipage}\par\addvspace{4pt}

\reserve{7\baselineskip}
% mdparagraph:P059:begin
\section*{12. Translation}
% mdparagraph:P059:end

\par\begin{minipage}{\linewidth}\raggedright
% item: NEW-U02-T01
% key: KEY-NEW-U02-T01
% mdparagraph:P060:begin
\itempara{1}{\textbf{Se mann sętte ðæt sǣd. Ðā męnn sęndon hit.} Singular subject \textbf{se mann} takes \textbf{sętte}; plural \textbf{ðā męnn} takes \textbf{sęndon}. Neuter accusative \textbf{hit} replaces \textbf{ðæt sǣd}. No direction phrase is needed to express “put down” in this assigned sense.}
% mdparagraph:P060:end

\end{minipage}\par\addvspace{4pt}

\par\begin{minipage}{\linewidth}\raggedright
% item: NEW-U02-T02
% key: KEY-NEW-U02-T02
% mdparagraph:P061:begin
\itempara{2}{\textbf{Ic mētte ðē. Ðū mēttest mē. Wit geręston unc.} The first two objects are accusative singular; \textbf{unc} is accusative dual and reflexive. The dual subject requires plural past \textbf{geręston}. Preserve \textbf{ē} in \textbf{mētte/\allowbreak{}mēttest} and the single \textbf{t} in \textbf{geręston}.}
% mdparagraph:P061:end

\end{minipage}\par\addvspace{4pt}

\par\begin{minipage}{\linewidth}\raggedright
% item: NEW-U02-T03
% key: KEY-NEW-U02-T03
% mdparagraph:P062:begin
\itempara{3}{\textbf{Hē nęrede ðone mann and lufode hine. Hē ne sęnde hine.} The first \textbf{hē} remains the subject of both \textbf{nęrede} and \textbf{lufode}. Both \textbf{hine} objects refer to the rescued man. Repeating \textbf{hē} before \textbf{lufode} is grammatical and preserves the meaning; a plural verb would not.}
% mdparagraph:P062:end

\end{minipage}\par\addvspace{4pt}

\par\begin{minipage}{\linewidth}\raggedright
% item: NEW-U02-T04
% key: KEY-NEW-U02-T04
% mdparagraph:P063:begin
\itempara{4}{\textbf{Ðū geblētsodest ðone dæg. Hīe gehālgodon hine.} The first verb is second singular preterite, the second third plural preterite. \textbf{Dæg} is masculine; the referring accusative is \textbf{hine}, naturally translated “it”. \textbf{Hit} would not preserve that grammatical gender.}
% mdparagraph:P063:end

\end{minipage}\par\addvspace{4pt}

\par\begin{minipage}{\linewidth}\raggedright
% item: NEW-U02-T05
% key: KEY-NEW-U02-T05
% mdparagraph:P064:begin
\itempara{5}{\textbf{“He made his dwelling. They performed their work. He rested himself in the dwelling.”} “Built” is suitable for \textbf{geworhte} here; “accomplished/\allowbreak{}completed” for \textbf{gefręmedon}; ordinary “rested” may omit English “himself”. \textbf{His/\allowbreak{}hiera} identify possessors, \textbf{hine} is reflexive accusative, and \textbf{on ðǣm hofe} gives a dative location. Do not translate \textbf{gefręmedon} as the separate unprefixed sense “benefited”.}
% mdparagraph:P064:end

\end{minipage}\par\addvspace{4pt}

\par\begin{minipage}{\linewidth}\raggedright
% item: NEW-U02-T06
% key: KEY-NEW-U02-T06
% mdparagraph:P065:begin
\itempara{6}{\textbf{Wē habbað ðæs mannes sǣd. Wē sęllað him wæter. Hē hīerð ūs.} \textbf{Ðæs mannes} is the limiting genitive before accusative \textbf{sǣd}; no extra head article is needed in the selected pattern. \textbf{Him} is dative recipient, \textbf{ūs} accusative object. The last verb is singular because its subject is \textbf{hē}.}
% mdparagraph:P065:end

\end{minipage}\par\addvspace{4pt}

\reserve{7\baselineskip}
% mdparagraph:P066:begin
\section*{13. Two connected readings}
% mdparagraph:P066:end

\par\begin{minipage}{\linewidth}\raggedright
% item: NEW-U02-N01
% key: KEY-NEW-U02-N01
% mdparagraph:P067:begin
\itempara{1}{A model translation prepared from the Old English:}
% mdparagraph:P067:end

% mdparagraph:P068:begin
“The man had seed and water. The men had a dwelling and birds. The birds lived in their dwelling. The man gave the men seed and water. They loved the birds and heard them. They gave them seed and water. They did not send the birds. They performed their work and rested themselves. The man said: ‘I love the birds.’”
% mdparagraph:P068:end

% mdparagraph:P069:begin
In the intended sequence the men receive seed and water, care for the birds, and give them seed and water. “House/\allowbreak{}home” for \textbf{hof}, “accomplished/\allowbreak{}completed” for \textbf{gefręmedon}, and plain “rested” are defensible. \textbf{Ic} in the final direct quotation is the speaking man; its present \textbf{lufie} is not changed to past just because \textbf{sægde} is past. A translation that consistently explains another plausible plural reference should be assessed for that interpretation, not rejected because it differs from the model’s unstated antecedents.
% mdparagraph:P069:end

\end{minipage}\par\addvspace{4pt}

\par\begin{minipage}{\linewidth}\raggedright
% item: NEW-U02-N02
% key: KEY-NEW-U02-N02
% mdparagraph:P070:begin
\itempara{2}{The intended references are:}
% mdparagraph:P070:end

% mdtable:T13:begin
\par\begingroup
\fontsize{10.5}{13.2}\selectfont
\setlength{\LTleft}{0pt}\setlength{\LTright}{\fill}
\noindent
\begin{tabular}{@{}>{\raggedright\arraybackslash}p{39mm}>{\raggedright\arraybackslash}p{45mm}>{\raggedright\arraybackslash}p{33mm}@{}}
\toprule
\textbf{Sentence and occurrence} & \textbf{Analysis} & \textbf{Intended referent} \\
\midrule
3: hiera & genitive plural possessor & the men \\
5: first Hīe & nominative plural & the men \\
5: second hīe & accusative plural & the birds \\
6: Hīe & nominative plural & the men \\
6: him & dative plural recipient & the birds \\
7: Hīe & nominative plural & the men \\
8: first Hīe & nominative plural & the men \\
8: hiera & genitive plural possessor & the men \\
8: final hīe & accusative plural, reflexive & the men \\
\bottomrule
\end{tabular}
\par\endgroup

% mdtable:T13:end
% mdparagraph:P071:begin
The subject verbs and the assigned constructions establish the cases of the \textbf{hīe} forms. They do not distinguish men from birds, both available plural antecedents. \textbf{Hiera} marks a plural possessor, not which plural group possesses. \textbf{Him} itself also permits masculine/\allowbreak{}neuter dative singular: an interpretation in which a suitable singular recipient is meant needs contextual explanation rather than a claim that the spelling is impossible. The model’s giving-and-care sequence favours men as givers and birds as recipients. All \textbf{ðā} in this account are articles, not the source reading’s temporal adverb.
% mdparagraph:P071:end

\end{minipage}\par\addvspace{4pt}

\par\begin{minipage}{\linewidth}\raggedright
% item: NEW-U02-N03
% key: KEY-NEW-U02-N03
% mdparagraph:P072:begin
\itempara{3}{\textbf{A to B:} “I have seed. You have water.” \textbf{Ic} is A; \textbf{ðū} is B.}
% mdparagraph:P072:end

% mdparagraph:P073:begin
\textbf{B to A:} “We two give the birds food. They hear us two.” \textbf{Wit} includes B and A; \textbf{hīe} refers to the birds, \textbf{unc} to A and B.
% mdparagraph:P073:end

% mdparagraph:P074:begin
\textbf{C to A and B:} “You two love the birds. I hear you two.” \textbf{Git} includes A and B; \textbf{ic} is C; \textbf{inc} refers to A and B.
% mdparagraph:P074:end

% mdparagraph:P075:begin
\textbf{A to B, C and D:} “We love them. You hear them, and they hear you.” By the given setting \textbf{wē} includes all four workers; \textbf{gē/\allowbreak{}ēow} identify B, C and D. The \textbf{hīe} forms refer to the birds. The speaker changes between turns; a first-person pronoun cannot be assigned permanently to A throughout the dialogue.
% mdparagraph:P075:end

\end{minipage}\par\addvspace{4pt}

\par\begin{minipage}{\linewidth}\raggedright
% item: NEW-U02-N04
% key: KEY-NEW-U02-N04
% mdparagraph:P076:begin
\itempara{4}{\textbf{Unc} is accusative dual, A and B, in \textbf{Hīe hīerað unc}: nominative plural \textbf{hīe}, the birds, controls \textbf{hīerað}. \textbf{Inc} is accusative dual, A and B, in \textbf{Ic hīere inc}: nominative singular \textbf{ic}, speaker C, controls \textbf{hīere}. \textbf{Ēow} is accusative plural, B, C and D, in \textbf{hīe hīerað ēow}: nominative plural \textbf{hīe}, the birds, controls \textbf{hīerað}. The object pronoun gives the people heard, not the finite verb’s person or number.}
% mdparagraph:P076:end

\end{minipage}\par\addvspace{4pt}

\par\begin{minipage}{\linewidth}\raggedright
% item: NEW-U02-N05
% key: KEY-NEW-U02-N05
% mdparagraph:P077:begin
\itempara{5}{\textbf{Se mann hæfð sǣd and wæter. Ðā męnn habbað hof and fuglas. Ðā fuglas libbað on hiera hofe. Se mann sęlð ðǣm mannum sǣd and wæter.}}
% mdparagraph:P077:end

% mdparagraph:P078:begin
The changes are \textbf{hæfde → hæfð}, \textbf{hæfdon → habbað}, \textbf{lifdon → libbað}, \textbf{sealde → sęlð}. Subjects, accusative objects, dative location and dative recipients remain unchanged. \textbf{Sęleð} is a valid source alternative to selected \textbf{sęlð}.
% mdparagraph:P078:end

\end{minipage}\par\addvspace{4pt}

\par\begin{minipage}{\linewidth}\raggedright
% item: NEW-U02-N06
% key: KEY-NEW-U02-N06
% mdparagraph:P079:begin
\itempara{6}{\textbf{“They performed their work and rested themselves.”} The first \textbf{hīe} is the nominative plural subject of both \textbf{gefręmedon} and \textbf{geręston}. Accusative \textbf{weorc} is the object of the first verb; the final accusative \textbf{hīe} is the reflexive object of the second. \textbf{Hiera} is a genitive plural possessor, referring to the men in the intended account. “Accomplished/\allowbreak{}completed their work and rested” is natural English. Do not analyse \textbf{hiera} as the subject or \textbf{hīe} as an independent new actor merely because it follows the second verb.}
% mdparagraph:P079:end

\end{minipage}\par\addvspace{4pt}

\reserve{7\baselineskip}
% mdparagraph:P080:begin
\section*{14. Source reading: the Sabbath}
% mdparagraph:P080:end

\par\begin{minipage}{\linewidth}\raggedright
% item: NEW-U02-R01
% key: KEY-NEW-U02-R01
% mdparagraph:P081:begin
\itempara{1}{A close model translation from the supplied Old English:}
% mdparagraph:P081:end

% mdparagraph:P082:begin
“Accordingly, then, the heavens and the earth and all their array were completed. And God then completed on the seventh day his work which he had made, and he rested himself on the seventh day from all the work which he had performed. And God blessed the seventh day and hallowed it, because on that day he ceased from his work which he created to make.”
% mdparagraph:P082:end

% mdparagraph:P083:begin
“Thus” may replace “accordingly”; “finished” may replace “completed”; “made holy/\allowbreak{}consecrated” may replace “hallowed”. English “rested” need not retain “himself”. Simple “made/\allowbreak{}performed” as well as “had made/\allowbreak{}had performed” can convey the relative past actions. Preserve the source’s distinct completing, resting, blessing and hallowing, and its final causal relation. \textbf{Hine} after hallowing refers to the day, not God.
% mdparagraph:P083:end

% mdparagraph:P084:begin
The ending deliberately remains close to \textbf{gescēop tō wyrceanne}. Cook’s own note gives “created to make”; a defensible close “created to do” may also convey the relation. The understood maker/\allowbreak{}doer is God. Neither “created and made” nor “finished creating” translates that relation as stated. The learner may consult the supplied explanation; independent mastery of this infinitive, the relative clauses or the passive is not being assessed. [1]
% mdparagraph:P084:end

\end{minipage}\par\addvspace{4pt}

\par\begin{minipage}{\linewidth}\raggedright
% item: NEW-U02-R02
% key: KEY-NEW-U02-R02
% mdparagraph:P085:begin
\itempara{2}{All six are \textbf{third-person singular preterite indicative}:}
% mdparagraph:P085:end

% mdtable:T14:begin
\par\begingroup
\fontsize{10.5}{13.2}\selectfont
\setlength{\LTleft}{0pt}\setlength{\LTright}{\fill}
\noindent
\begin{tabular}{@{}>{\raggedright\arraybackslash}p{31mm}>{\raggedright\arraybackslash}p{28mm}>{\raggedright\arraybackslash}p{24mm}>{\raggedright\arraybackslash}p{32mm}@{}}
\toprule
\textbf{Form in the passage} & \textbf{Headword} & \textbf{Class} & \textbf{Meaning here} \\
\midrule
gefylde & gefyllan & I & completed/\allowbreak{}finished \\
geworhte & gewyrcean & I & made \\
geręste & geręstan & I & rested \\
gefręmede & gefręmman & I & performed/\allowbreak{}accomplished \\
geblētsode & geblētsian & II & blessed \\
gehālgode & gehālgian & II & hallowed/\allowbreak{}made holy \\
\bottomrule
\end{tabular}
\par\endgroup

% mdtable:T14:end
% mdparagraph:P086:begin
The singular subject fixes third person here, although these verb spellings can be first singular elsewhere. Do not count the supplied strong \textbf{geswāc/\allowbreak{}gescēop}, familiar \textbf{wǣron}, the participle \textbf{fullfręmede}, or infinitive \textbf{wyrceanne} as one of the six requested weak preterites. Headword lookup is expressly allowed.
% mdparagraph:P086:end

\end{minipage}\par\addvspace{4pt}

\par\begin{minipage}{\linewidth}\raggedright
% item: NEW-U02-R03
% key: KEY-NEW-U02-R03
% mdparagraph:P087:begin
\itempara{3}{In \textbf{geręste hine}, \textbf{hine} is masculine accusative singular referring back to \textbf{God/\allowbreak{}hē}, the subject: “rested himself”, naturally “rested”. It is reflexive. In \textbf{hine gehālgode}, \textbf{hine} is masculine accusative singular referring to \textbf{ðone seofoðan dæg}: “hallowed it/\allowbreak{}made it holy”. It is an ordinary object, not reflexive. The identical pronoun form expresses different reference relations; it does not inherently mean “self”.}
% mdparagraph:P087:end

\end{minipage}\par\addvspace{4pt}

\par\begin{minipage}{\linewidth}\raggedright
% item: NEW-U02-R04
% key: KEY-NEW-U02-R04
% mdparagraph:P088:begin
\itempara{4}{\textbf{Hiera frætwung}: the possessors are collectively the heavens and earth; the possessed noun is nominative feminine singular \textbf{frætwung}, their array. \textbf{His weorc}: the possessor is God; the possessed noun is accusative neuter singular \textbf{weorc}, his work. \textbf{His weorces}: the same possessor, God, with genitive neuter singular \textbf{weorces}, governed by supplied \textbf{geswāc}. \textbf{Hiera/\allowbreak{}his} are genitive pronouns selected by the possessor; they do not inflect with the possessed head. In these actual phrases all three possessed heads are singular; the two work forms specifically illustrate a change of case. A change of head number would likewise leave the third-person possessor unchanged.}
% mdparagraph:P088:end

\end{minipage}\par\addvspace{4pt}

\par\begin{minipage}{\linewidth}\raggedright
% item: NEW-U02-R05
% key: KEY-NEW-U02-R05
% mdparagraph:P089:begin
\itempara{5}{\textbf{Heofonas}: nominative plural. \textbf{Eorðe}: nominative singular. \textbf{Frætwung}: nominative singular. The complete coordinated subject is \textbf{heofonas and eorðe and eall hiera frætwung}, “the heavens and the earth and all their array”. It is the subject of plural \textbf{wǣron} in the supplied passive expression. \textbf{Hiera} is a possessor inside the last noun phrase; \textbf{fullfręmede} is the supplied participle, not a fourth subject.}
% mdparagraph:P089:end

\end{minipage}\par\addvspace{4pt}

\par\begin{minipage}{\linewidth}\raggedright
% item: NEW-U02-R06
% key: KEY-NEW-U02-R06
% mdparagraph:P090:begin
\itempara{6}{God \textbf{completed his work}, \textbf{rested from it}, \textbf{blessed the seventh day}, and \textbf{hallowed that day}. The causal clause also describes his \textbf{ceasing from his work on that day}, using supplied \textbf{geswāc}. The temporal phrases are supplied as “on the seventh day/\allowbreak{}on that day”. The direct object of \textbf{geblētsode} is the day; the following \textbf{hine} resumes that day. An answer that makes God hallow himself changes the source’s reference.}
% mdparagraph:P090:end

\end{minipage}\par\addvspace{4pt}

\par\begin{minipage}{\linewidth}\raggedright
% item: NEW-U02-R07
% key: KEY-NEW-U02-R07
% mdparagraph:P091:begin
\itempara{7}{The reason is that \textbf{on that day God ceased from his work}. The following relative phrase describes that work as what he \textbf{created to make/\allowbreak{}do}. The genitive after “ceased”, the causal link, and the entire relative/\allowbreak{}infinitive phrase come from the supplied help. God remains the understood maker; the work is not made into a new subject of \textbf{wyrceanne}. Preserve this relation without requiring one uniquely smooth English paraphrase. [1]}
% mdparagraph:P091:end

\end{minipage}\par\addvspace{4pt}

\reserve{7\baselineskip}
% mdparagraph:P092:begin
\section*{15. Review}
% mdparagraph:P092:end

\par\begin{minipage}{\linewidth}\raggedright
% item: NEW-U02-X01
% key: KEY-NEW-U02-X01
% mdparagraph:P093:begin
\itempara{1}{\textbf{Wit hæfdon sǣd. Wē sealdon ðǣm fuglum męte. Hīe ne hīerdon ūs.} “We two had seed. We gave the birds food. They did not hear us.” All three finite verbs become preterite; only the last clause receives \textbf{ne}. Dual \textbf{wit} still takes a plural form, and the noun/\allowbreak{}pronoun cases remain unchanged.}
% mdparagraph:P093:end

\end{minipage}\par\addvspace{4pt}

\par\begin{minipage}{\linewidth}\raggedright
% item: NEW-U02-X02
% key: KEY-NEW-U02-X02
% mdparagraph:P094:begin
\itempara{2}{\textbf{Hē sęlð mē sǣd. Wit habbað hit. Ðā fuglas hīerað unc.} The four repairs are \textbf{Hine → Hē} for a nominative subject; \textbf{ic → mē} for a dative recipient; \textbf{hæfst → habbað} for agreement with dual \textbf{wit}; \textbf{ūre → unc} for an accusative dual object. The rejected forms are real forms used in the wrong roles, not invented spellings. \textbf{Hit} correctly refers to neuter \textbf{sǣd}.}
% mdparagraph:P094:end

\end{minipage}\par\addvspace{4pt}

\par\begin{minipage}{\linewidth}\raggedright
% item: NEW-U02-X03
% key: KEY-NEW-U02-X03
% mdparagraph:P095:begin
\itempara{3}{\textbf{Se mann hæfde fuglas. Hē sealde him sǣd. Hē geręste hine.} “The man had birds. He gave them seed. He rested himself.” \textbf{Hē} refers to the introduced man; dative plural \textbf{him} to the birds; accusative \textbf{hine} reflexively to the man. Accept grammatical word-order or documented spelling alternatives that preserve these past actions and references. Replacing every later pronoun with the noun could express similar content but would not satisfy the requested pronoun practice.}
% mdparagraph:P095:end

\end{minipage}\par\addvspace{4pt}

\clearpage
% mdparagraph:P096:begin
\section*{Source notes}
% mdparagraph:P096:end

\begingroup\fontsize{10}{12.8}\selectfont
% mdparagraph:P097:begin
[1] Cook, \emph{A First Book in Old English}, second edition (1895), Reader I, p.128 /\allowbreak{} one-based PDF144, final paragraph, editorial prose lines 1–7 (only marginal 5 printed). The model follows the checked Old English rather than a modern biblical version. Note 12 supplies \textbf{creavit ut faceret}, “created to make”; §200, p.106 /\allowbreak{} PDF122, explains the inflected infinitive. The temporal accusative, cessation genitive and relative-clause support follow the precise locators in the learner’s source notes.
% mdparagraph:P097:end

% mdparagraph:P098:begin
The finite and pronoun models use Cook §§81–83,96,113–124 and the audited lexical entries. Selected contractions and alternatives were independently checked against Wright and Wright, \emph{Old English Grammar}, second edition (1914), §§305,476,524,528–529,532–533. Dual agreement is explicitly supported by Ringe and Taylor, \emph{The Development of Old English} (2014), §1.1.2, p.6. The first unit’s noun/\allowbreak{}article paradigms remain the basis of the nominal forms; no inherited-course milestone is assumed.
% mdparagraph:P098:end

\endgroup
\end{document}
